%% manuscript_pt.tex — Versão em português para leitura
%% Tradução do manuscrito submetido à Lancet Regional Health – Americas
%% ===================================================================
\documentclass[12pt]{article}

\usepackage[utf8]{inputenc}
\usepackage[T1]{fontenc}
\usepackage[brazilian]{babel}
\usepackage{amsmath,amssymb}
\usepackage{graphicx}
\usepackage{booktabs}
\usepackage{multirow}
\usepackage{xcolor}
\usepackage{hyperref}
\usepackage[margin=2.5cm]{geometry}
\usepackage{setspace}
\onehalfspacing

\begin{document}

\begin{center}
{\LARGE\bfseries Efeitos não lineares e defasados dos extremos de frio
nas internações hospitalares respiratórias em uma metrópole subtropical:
análise de séries temporais quasi-Poisson em Curitiba, Brasil}\\[1.5em]
{\large Felipe Baglioli$^{1,*}$, Lucas Rover$^1$}\\[0.5em]
$^1$Universidade Federal do Paraná (UFPR), Curitiba, PR, Brasil\\
$^*$Autor correspondente: felipe.baglioli@ufpr.br\\[1em]
{\small\itshape Versão em português para leitura --- O manuscrito de
submissão está em inglês}
\end{center}

\bigskip
\hrule
\bigskip

%% ═══════════════════════════════════════════════════════════════════
%% RESUMO
%% ═══════════════════════════════════════════════════════════════════

\section*{Resumo}

\textbf{Contexto}\quad
Os efeitos dos extremos de temperatura na saúde permanecem pouco
caracterizados em cidades subtropicais do Hemisfério Sul, onde ondas
de frio produzem perfis de exposição distintos daqueles estudados em
climas temperados e tropicais.

\textbf{Métodos}\quad
Analisamos a temperatura mínima diária (INMET, 1961--2024) e as
internações hospitalares respiratórias (DATASUS SIH/AIH, CID-10
J00--J99, 2022--2024; $N$=42.716) em Curitiba, Brasil
($\sim$1,8~milhão de habitantes). Modelos distribuídos de defasagem
não linear (DLNM) com regressão quasi-Poisson quantificaram
associações exposição--resposta não lineares e defasadas ao longo de
0--21~dias, ajustando por dia da semana, feriados e tendência de longo
prazo via \textit{spline} natural (7~gl/ano). A temperatura de mínima
morbidade (TMM) foi identificada por busca em grade. Análises de
sensibilidade variaram janelas de defasagem, graus de liberdade
sazonais e excluíram 2022 por efeitos residuais da COVID-19.

\textbf{Resultados}\quad
Em 1.096~dias de estudo, foram registradas 42.716~internações
respiratórias (51\% masculinos; 38\% com $\geq$65~anos) e
3.715~óbitos intra-hospitalares. O DLNM identificou TMM=17,8\,°C.
Extremos de frio no 1º percentil de $T_\mathrm{min}$ (5,2\,°C)
apresentaram RR cumulativo=1,42 [IC 95\%: 1,09--1,84] ao longo de
0--21~dias; no 5º percentil (8,5\,°C), RR=1,22 [1,02--1,44]. A
estrutura de defasagem teve pico entre 3--10~dias. Efeitos do calor
não foram significativos (P90 RR=1,06 [0,98--1,14]). Os resultados
foram robustos nas análises de sensibilidade
($\hat{\phi}$=1,32; Ljung--Box $p$=0,14).

\textbf{Interpretação}\quad
Extremos de frio, e não de calor, são o principal fator de internações
hospitalares respiratórias neste contexto subtropical, com efeitos
defasados que atingem o pico entre 3--10~dias. Sistemas de alerta
precoce em saúde pública devem priorizar previsões de ondas de frio
junto aos alertas de calor.

\textbf{Financiamento}\quad
Nenhum.

\bigskip

\textbf{Palavras-chave:} extremos de temperatura; ondas de frio;
internações hospitalares; modelo distribuído de defasagem não linear;
clima e saúde; morbidade respiratória; subtropical; Curitiba

\bigskip
\hrule
\bigskip

%% ═══════════════════════════════════════════════════════════════════
%% PESQUISA EM CONTEXTO
%% ═══════════════════════════════════════════════════════════════════

\section*{Pesquisa em Contexto}

\textbf{Evidências antes deste estudo}

Pesquisamos no PubMed e Scopus estudos publicados entre janeiro de
2010 e dezembro de 2025 usando os termos ``extremos de temperatura'',
``internações hospitalares'', ``defasagem distribuída'' e ``Hemisfério
Sul'' ou ``América Latina''. Grandes estudos multicidades usando
métodos DLNM estabeleceram curvas temperatura--mortalidade em forma de
U globalmente, mas cidades subtropicais no Hemisfério Sul ---
particularmente na América do Sul --- permanecem sub-representadas.
Poucos estudos quantificaram os efeitos não lineares e defasados dos
extremos de frio na morbidade respiratória nesta região.

\textbf{Valor agregado deste estudo}

Este estudo fornece a primeira análise DLNM de extremos de temperatura
e internações hospitalares respiratórias em uma cidade subtropical
sul-americana. Extremos de frio no 1º percentil de $T_\mathrm{min}$
estão associados a um aumento de 42\% nas internações respiratórias
(RR=1,42 [1,09--1,84]) com pico defasado entre 3--10~dias --- uma
estrutura de defasagem não documentada anteriormente com métodos
formais de DLNM nesta região.

\textbf{Implicações de todas as evidências disponíveis}

Sistemas de alerta precoce em saúde pública em cidades subtropicais
devem incorporar previsões de ondas de frio junto aos alertas de
calor, em vez de focar exclusivamente em extremos de calor. A
estrutura de defasagem (3--10~dias) sugere uma janela para intervenção
preventiva após o início da onda de frio.

%% ═══════════════════════════════════════════════════════════════════
%% INTRODUÇÃO
%% ═══════════════════════════════════════════════════════════════════

\section{Introdução}

A relação entre extremos de temperatura e saúde humana é reconhecida
há muito tempo, com aumentos documentados na mortalidade durante ondas
de calor em diferentes regiões. Certos grupos --- idosos, crianças e
pessoas com condições pré-existentes --- são desproporcionalmente
afetados.

Extremos de frio representam uma ameaça comparável, e em alguns
contextos maior, à saúde pública. Temperaturas ambientais baixas
desencadeiam vasoconstrição, elevam a pressão arterial e aumentam a
viscosidade sanguínea, elevando o risco cardiovascular. Infecções
respiratórias atingem o pico durante períodos frios devido à redução
da depuração mucociliar e ao aumento do confinamento em ambientes
fechados. Diferentemente da morbidade relacionada ao calor, que se
manifesta em 0--3~dias, os efeitos relacionados ao frio exibem
defasagens de 3--14~dias, complicando a atribuição em estudos
ecológicos.

No contexto das mudanças climáticas, atividades antropogênicas
aumentaram a frequência e intensidade de eventos extremos e seus
impactos na saúde. Poluentes atmosféricos afetam os sistemas
respiratório, cardiovascular e neural, e episódios extremos de
poluição permanecem associados a internações hospitalares e
mortalidade elevadas.

Apesar dessas evidências, poucos estudos quantificaram os efeitos não
lineares e defasados dos extremos de temperatura na morbidade
respiratória em cidades subtropicais do Hemisfério Sul usando métodos
inferenciais formais. Curitiba, metrópole subtropical no sul do Brasil
($\sim$930~m de altitude, clima Cfb), oferece um estudo de caso
relevante devido às suas temperaturas de inverno abaixo de 5\,°C e
episódios sazonais de poluição. Este estudo utiliza modelos
distribuídos de defasagem não linear para quantificar a relação
exposição--resposta entre temperatura mínima diária e internações
hospitalares respiratórias em Curitiba (2022--2024), com análises
suplementares de agrupamento e aprendizado de máquina.

%% ═══════════════════════════════════════════════════════════════════
%% MÉTODOS
%% ═══════════════════════════════════════════════════════════════════

\section{Métodos}

\subsection{Desenho do estudo e cenário}

Trata-se de um estudo ecológico de séries temporais examinando a
associação entre exposições ambientais diárias e internações
hospitalares respiratórias em Curitiba, Paraná, Brasil
(população~$\sim$1,8~milhão), de 1º de janeiro de 2022 a 31 de
dezembro de 2024 (1.096~dias). O clima subtropical de Curitiba
apresenta temperaturas de inverno abaixo de 5\,°C e máximas de verão
próximas de 30\,°C.

\subsection{Fontes de dados}

Registros individuais de internação hospitalar foram obtidos do
Sistema Unificado de Saúde (SIH/SUS--DATASUS) para 2022--2024
($N$=42.716~internações respiratórias, CID-10 J00--J99). O desfecho
primário foi a contagem diária de internações respiratórias.
Demografia: 51\% masculinos, 38\% com $\geq$65~anos. Mortalidade
intra-hospitalar: 8,7\% ($n$=3.715).

Temperatura mínima ($T_\mathrm{min}$), média e máxima diárias,
umidade relativa e pressão atmosférica foram obtidas da estação
automática INMET~A807. Para cálculos de percentis ETCCDI, o registro
completo (1961--2024) foi utilizado; análises combinadas foram
restritas a 2022--2024.

Concentrações de PM$_{2,5}$ foram medidas com sensor PurpleAir
Flex-II (Sensor ID~90875) em resolução diária. Leituras brutas foram
corrigidas pela equação de correção EPA:
%
\begin{equation}
  \mathrm{PM}_{2{,}5}^{\mathrm{corr}} =
    0{,}524 \times \mathrm{PM}_{2{,}5}^{\mathrm{sensor}}
    - 0{,}0862 \times \mathrm{UR} + 5{,}75
\end{equation}
%
Esta correção foi validada contra monitores de Método Equivalente
Federal e é amplamente adotada para dados PurpleAir. Um protocolo de
controle de qualidade intercanal excluiu dias com diferenças
intercanais superiores a 5~µg/m$^3$ e 70\% de diferença relativa
simultaneamente. Após QC, 697~dias (63,6\%) possuíam PM$_{2,5}$
válido corrigido pela EPA.

\subsection{Definições e dados faltantes}

Dias de temperatura extrema foram definidos usando percentis mensais
ETCCDI (TX90p para extremamente quente, TN10p para extremamente frio)
sobre a linha de base 1961--2024. Ondas de calor:
$\geq$3~dias consecutivos extremamente quentes; ondas de frio:
$\geq$3~dias consecutivos extremamente frios. Taxas de dados
faltantes: $T_\mathrm{min}$~2,6\%, PM$_{2,5}$~36,4\%,
internações~0,2\%. Análises primárias usaram casos completos; análises
de sensibilidade com imputação produziram resultados similares
(Suplemento).

\subsection{Análise estatística}

Séries temporais foram decompostas usando STL para identificar
componentes de tendência e sazonalidade. Testes de Dickey--Fuller
aumentado (ADF) avaliaram estacionaridade. ACP foi realizada sobre
PM$_{2,5}$, $T_\mathrm{min}$ e internações padronizados ($N$=683~dias
de casos completos). Agrupamento K-médias foi aplicado a seis
variáveis padronizadas, com $K$ ótimo selecionado pelo índice de
silhueta. Funções de correlação cruzada (defasagens 0--28~dias) foram
calculadas entre exposições e internações.

\subsection{Previsão por aprendizado de máquina}

Como análise suplementar, três modelos --- XGBoost, perceptron
multicamada e máquina de aprendizado extremo --- foram treinados
para prever internações diárias nas defasagens 0--7~dias. Valores
SHAP quantificaram a importância das variáveis. Detalhes, resultados
e projeções CMIP6 são relatados no Suplemento.

\subsection{Modelo distribuído de defasagem não linear}

Ajustamos modelos distribuídos de defasagem não linear (DLNM) para
quantificar relações exposição--resposta não lineares e defasadas:
%
\begin{equation}
  \log\!\big(\mathrm{E}[Y_t]\big) = \alpha
    + \mathbf{s}\bigl(x_t, \mathrm{lag}; \boldsymbol{\eta}\bigr)
    + \sum_{d=1}^{6} \beta_d \,\mathrm{DS}_{d,t}
    + \gamma \,\mathrm{Feriado}_t
    + \mathrm{ns}(\mathrm{tempo}_t,\, 7\,\mathrm{gl/ano})
\end{equation}
%
onde $Y_t$ é a contagem diária de internações respiratórias
(quasi-Poisson), $\mathbf{s}(\cdot)$ é a função de base cruzada
(\textit{splines} naturais cúbicos: gl exposição=4, gl defasagem=5),
DS denota indicadores de dia da semana e
$\mathrm{ns}(\mathrm{tempo}, 21\,\mathrm{gl})$ captura tendência
e sazonalidade. A defasagem máxima foi de 21~dias. A temperatura de
mínima morbidade (TMM) foi identificada por busca em grade, seguindo
Gasparrini et al. Intervalos de confiança foram calculados pelo método
delta. Análises de sensibilidade variaram a defasagem máxima (14, 21,
28~dias), gl sazonais (6, 7, 8~gl/ano) e excluíram 2022 para avaliar
efeitos residuais da COVID-19.

\paragraph{Ética e reporte}
Este estudo utilizou exclusivamente dados agregados e publicamente
disponíveis do DATASUS, INMET e PurpleAir. Conforme a Resolução
CNS~510/2016, pesquisas com dados públicos agregados são isentas de
apreciação por comitê de ética. Este estudo segue as diretrizes
STROBE.

%% ═══════════════════════════════════════════════════════════════════
%% RESULTADOS
%% ═══════════════════════════════════════════════════════════════════

\section{Resultados}

\subsection{Estatísticas descritivas}

A Tabela~\ref{tab:descritiva} apresenta estatísticas descritivas para
todas as variáveis. Um total de 42.716~internações respiratórias e
3.715~óbitos foram registrados em 1.096~dias. As internações diárias
médias foram de 32,3 (DP~11,7; variação 1--73). O PM$_{2,5}$
corrigido pela EPA teve média de 5,84~µg/m$^3$ (mediana~4,12), com
4,7\% dos dias excedendo a diretriz de 24~horas da OMS
(15~µg/m$^3$). Identificamos 215~dias extremamente quentes (20,1\%) e
39~dias extremamente frios (3,7\%), agregando-se em 25~ondas de calor
e 4~ondas de frio (eTabela~1).

\begin{table}[ht]
\centering
\caption{Estatísticas descritivas para todas as variáveis do estudo,
Curitiba, 2022--2024 ($N$=1.096~dias).}
\label{tab:descritiva}
\begin{tabular}{lrrrrrrrr}
\toprule
Variável & $N$ & Média & DP & Mediana & Q25 & Q75 & Mín & Máx \\
\midrule
$T_\mathrm{máx}$ (°C) & 1070 & 24,53 & 4,75 & 24,80 & 21,50 & 28,10 & 10,20 & 35,40 \\
$T_\mathrm{méd}$ (°C) & 1042 & 18,74 & 3,63 & 18,90 & 16,30 & 21,50 &  7,80 & 27,30 \\
$T_\mathrm{mín}$ (°C) & 1068 & 14,75 & 3,75 & 15,20 & 12,00 & 17,70 &  1,30 & 23,40 \\
PM$_{2,5}$ (µg/m$^3$) &  697 &  5,84 & 7,40 &  4,12 &  2,43 &  6,55 &  0,00 & 65,19 \\
Umid.\ rel.\ (\%)     &  859 & 79,98 &  9,52 & 80,90 & 74,15 & 86,55 & 43,00 & 98,40 \\
Intern.\ (n/dia)      & 1094 & 32,29 & 11,74 & 32,00 & 23,00 & 41,00 & 1,00 & 73,00 \\
\bottomrule
\end{tabular}
\end{table}

\subsection{Análises exploratórias}

Análises de agrupamento K-médias e correlação defasada são relatadas
no Suplemento (eTabelas~2--5; eFiguras~1--4). Resumidamente, três
agrupamentos clima--poluição emergiram ($K$=3, silhueta=0,357), com
o agrupamento Frio Limpo mostrando correlações positivas defasadas em
3--10~dias. A análise de defasagem estendida revelou correlação máxima
$T_\mathrm{mín}$--internação na defasagem~9 ($r$=$-$0,320;
$p<$0,001), motivando a janela de defasagem DLNM de 0--21~dias.

\subsection{Modelos distribuídos de defasagem não linear}

\paragraph{Exposição--resposta de temperatura}
O DLNM quasi-Poisson para $T_\mathrm{mín}$ ($N$=1.068;
$\hat{\phi}$=1,32) identificou TMM=17,8\,°C
(Figura~\ref{fig:dlnm_temp_geral}). Extremos de frio mostraram RR
cumulativo monotonicamente crescente abaixo da TMM: no 1º percentil
(5,2\,°C), RR=1,42 [1,09--1,84]; no 5º percentil (8,5\,°C),
RR=1,22 [1,02--1,44]; no 10º percentil (9,8\,°C), RR=1,17
[1,00--1,37] (Tabela~\ref{tab:rr_temp}). Efeitos do calor não foram
significativos no 90º percentil (RR=1,06 [0,98--1,14]). A estrutura
de defasagem no frio (P10) teve pico nas defasagens 3--10~dias
(Figura~\ref{fig:dlnm_temp_lag}a), consistente com os resultados de
correlação cruzada.

\paragraph{Sensibilidade e diagnósticos}
Os resultados foram robustos à janela de defasagem e gl sazonais
(Tabela~\ref{tab:sensibilidade}). O RR cumulativo no P10 variou de
1,05 (lag=14) a 1,23 (lag=28). Excluir 2022 fortaleceu o efeito
(RR=1,37 [1,11--1,70]). O teste ADF nas internações brutas confirmou
não estacionaridade ($p$=0,164); resíduos STL foram estacionários
($p$=0,0003; eFigura~7). Diagnósticos do modelo não mostraram
autocorrelação (Durbin--Watson=2,00; Ljung--Box $p$=0,14; eTabela~7).

\begin{table}[ht]
\centering
\caption{Risco relativo (RR) cumulativo em percentis-chave de
$T_\mathrm{mín}$, centrado na TMM=17,8\,°C. DLNM quasi-Poisson,
defasagem 0--21~dias, $N$=1.068.}
\label{tab:rr_temp}
\begin{tabular}{lrrr}
\toprule
Percentil & $T_\mathrm{mín}$ (°C) & RR & IC 95\% \\
\midrule
P1  &  5,2 & 1,418 & 1,091--1,842 \\
P5  &  8,5 & 1,215 & 1,023--1,442 \\
P10 &  9,8 & 1,167 & 0,998--1,365 \\
P25 & 12,0 & 1,133 & 0,979--1,312 \\
P75 & 17,7 & 1,001 & 0,996--1,005 \\
P90 & 19,4 & 1,055 & 0,978--1,138 \\
P95 & 20,0 & 1,102 & 0,979--1,240 \\
P99 & 21,4 & 1,285 & 0,983--1,679 \\
\bottomrule
\end{tabular}
\end{table}

\begin{table}[ht]
\centering
\caption{Análises de sensibilidade para o DLNM de temperatura. RR
cumulativo no 10º percentil de $T_\mathrm{mín}$ (9,8\,°C).}
\label{tab:sensibilidade}
\begin{tabular}{lrrrr}
\toprule
Cenário & RR & IC 95\% & $\hat{\phi}$ & $N$ \\
\midrule
Principal (lag=21, 7\,gl/ano) & 1,17 & 1,00--1,37 & 1,32 & 1068 \\
Lag=14                        & 1,05 & 0,92--1,18 & 1,31 & 1068 \\
Lag=28                        & 1,23 & 1,01--1,51 & 1,32 & 1068 \\
Sazonal 6\,gl/ano             & 1,17 & 1,00--1,36 & 1,33 & 1068 \\
Sazonal 8\,gl/ano             & 1,21 & 1,03--1,42 & 1,31 & 1068 \\
Excluir 2022                  & 1,37 & 1,11--1,70 & 1,26 &  718 \\
\bottomrule
\end{tabular}
\end{table}

\begin{figure}[ht]
\centering
\includegraphics[width=0.85\textwidth]{../analysis/dlnm/figures/fig_dlnm_temp_overall.pdf}
\caption{Curva exposição--resposta cumulativa para $T_\mathrm{mín}$
e internações respiratórias (DLNM, defasagem 0--21~dias). Linha
tracejada vermelha: TMM (17,8\,°C). Área sombreada: IC 95\%.}
\label{fig:dlnm_temp_geral}
\end{figure}

\begin{figure}[ht]
\centering
\includegraphics[width=\textwidth]{../analysis/dlnm/figures/fig_dlnm_temp_lag.pdf}
\caption{RR específico por defasagem para (a) extremo de frio (P10,
9,8\,°C) e (b) extremo de calor (P90, 19,4\,°C), centrado na
TMM=17,8\,°C.}
\label{fig:dlnm_temp_lag}
\end{figure}

\subsection{Previsão por aprendizado de máquina}

Análises suplementares de ML (eTabela~8, eFigura~8) mostraram o
XGBoost como melhor modelo (MAPE=19,4\% na defasagem~0;
RMSE=8,37), com a temperatura classificada entre os principais
preditores pela análise SHAP.

%% ═══════════════════════════════════════════════════════════════════
%% DISCUSSÃO
%% ═══════════════════════════════════════════════════════════════════

\section{Discussão}

Nossa análise DLNM fornece evidência inferencial formal de que
extremos de frio são o principal fator relacionado à temperatura das
internações respiratórias em Curitiba. No 1º percentil de
$T_\mathrm{mín}$ (5,2\,°C), o RR cumulativo atingiu 1,42
[1,09--1,84] ao longo de 0--21~dias, com pico de defasagem entre
3--10~dias. Essa estrutura de defasagem é consistente com a
fisiopatologia conhecida do frio: vasoconstrição induzida pelo frio,
viscosidade sanguínea elevada e maior suscetibilidade a infecções
respiratórias se manifestam ao longo de dias, não horas.

Esses achados podem ser comparados com o estudo multicountry de
Gasparrini et al., que reportou curvas temperatura--mortalidade em
forma de U em 384~localidades usando o mesmo arcabouço DLNM. Nossa
TMM de 17,8\,°C é consistente com TMMs subtropicais reportadas naquele
estudo, e o risco atribuível ao frio no P1 é comparável em magnitude
ao excesso de mortalidade relacionado ao frio em climas temperados.
Notavelmente, Gasparrini et al. não incluíram Curitiba ou cidades
comparáveis do sul do Brasil, e nosso estudo preenche essa lacuna
fornecendo as primeiras estimativas de morbidade baseadas em DLNM para
este contexto subtropical. Contudo, como estudo de cidade única, nossos
resultados carecem da força estatística do \textit{pooling} multicidade.

As análises de sensibilidade apoiam a robustez do efeito do frio. A
direção e magnitude foram consistentes entre janelas de defasagem
(14--28~dias) e gl sazonais (6--8/ano). Excluir 2022 fortaleceu o
efeito (RR=1,37 [1,11--1,70]), sugerindo que efeitos residuais da
COVID-19 podem atenuar a associação temperatura--morbidade no período
completo.

Esses resultados têm implicações para sistemas de alerta precoce em
saúde pública em contextos subtropicais. Previsões de ondas de frio
devem ser priorizadas junto aos alertas de calor, particularmente dada
a estrutura de defasagem (3--10~dias) que oferece uma janela para
intervenção preventiva.

\subsection{Limitações}

Este estudo ecológico de séries temporais não permite inferência causal
em nível individual. Não havia dados individuais de exposição,
comportamento ou comorbidades. O período de análise de 3 anos
(2022--2024) limita o poder para eventos raros (ondas de frio, $n$=4)
e fornece apenas três realizações do ciclo sazonal. Confundimento por
circulação viral respiratória (influenza, VSR) não pode ser excluído,
pois dados de vigilância viral não foram incluídos. Contagens totais
diárias foram usadas sem estratificação por idade ou subgrupo CID-10;
análises por causa específica são prioridade para trabalhos futuros.
Os dados de temperatura são de uma única estação INMET (A807), que
pode não capturar variação microclimática na área metropolitana. Como
estudo de cidade única, os resultados podem não ser generalizáveis;
análises multicidade \textit{pooled} seguindo Gasparrini et al.
forneceriam evidências mais robustas.

\subsection{Conclusão}

Este estudo fornece evidências de que, em uma metrópole subtropical,
extremos de frio estão associados ao aumento das internações
hospitalares respiratórias, com estrutura de defasagem atingindo pico
entre 3--10~dias. O DLNM identificou TMM=17,8\,°C e estimou RR
cumulativo=1,42 [1,09--1,84] no frio extremo (P1, 5,2\,°C) ao longo
de 0--21~dias. Efeitos do calor não foram significativos. Sistemas de
saúde pública em cidades subtropicais do Hemisfério Sul devem integrar
previsões de ondas de frio em sistemas de alerta precoce. Trabalhos
futuros devem estender esta análise a múltiplas cidades com desfechos
específicos por causa e estratificados por idade.

%% ═══════════════════════════════════════════════════════════════════
%% DECLARAÇÕES
%% ═══════════════════════════════════════════════════════════════════

\section*{Contribuições}
FB: Conceitualização, Curadoria de dados, Análise formal, Investigação,
Metodologia, Redação --- rascunho original. LR: Metodologia, Software,
Validação, Análise formal, Redação --- revisão e edição. Todos os
autores tiveram acesso completo a todos os dados do estudo e tiveram
responsabilidade final pela decisão de submeter para publicação.

\section*{Declaração de interesses}
Os autores declaram não haver interesses concorrentes.

\section*{Compartilhamento de dados}
Todo código e dados processados estão disponíveis em
\url{https://github.com/Roverlucas/cwb-temp-pollution-health}.
Dados brutos de internação: DATASUS
(\url{https://datasus.saude.gov.br/}); dados meteorológicos: INMET
estação~A807; PM$_{2,5}$: PurpleAir Sensor ID~90875.

\section*{Financiamento}
Esta pesquisa não recebeu nenhum financiamento específico de agências
nos setores público, comercial ou sem fins lucrativos.

\section*{Agradecimentos}
Os autores agradecem ao Instituto Nacional de Meteorologia (INMET),
DATASUS e PurpleAir Inc.\ por fornecerem dados abertos.

\end{document}
